\documentclass{report}
\usepackage{amsmath,amssymb}
\setlength{\parindent}{0mm}
\setlength{\parskip}{1em}
\begin{document}
\begin{center}
\rule{6in}{1pt} \
{\large Denis Jarema \\
{\bf A Multiscale Approach for Particle Transport Simulation in Low Reynolds Number Flows}}

Traunsteiner Str 5/0405 \\ 81549 Munich \\ Germany
\\
{\tt denis.jarema@mytum.de}\\
Philipp Neumann\\
Tobias Weinzierl\end{center}

In this paper, we study a nanoparticle suspended in water within a
channel with a complex wall geometry in a laminar flow regime. We combine
two independent iterative methods into one multiscale approach to compute
the movement of the particle: a Lattice Boltzmann fluid-structure
interaction code and a Navier-Stokes solver with the Fax\'{e}n particle
force estimation. A new switching strategy between these two methods is
developed. On the one hand, the computationally expensive Lattice
Boltzmann fluid-structure code is used in time intervals when short-time
effects may yield a strong impact on the particle simulation. On the
other hand, the Navier-Stokes solver with the Fax\'{e}n correction is
used to compute the particle movement when long-time predictions to the
particle motion are possible and sufficient. We describe our coupling
strategy, the mapping of unknowns between the two solvers and provide
results for different particulate flow scenarios. Due to coupling of the
two systems and the automatic switching we reduce the total computing
time by magnitudes.


\end{document}
